Of the test articles, \val{ts.origin.publisher_section_jsonld.pct} carry JSON-LD section labels and \val{ts.origin.publisher_category.pct} publisher categories. These labels are editorial topics. Section conventions differ between publishers, and one section reflects one editorial view of an article that may cover several topics. The mapping table is conservative, validated against the IPTC tree, and reviewed by the first author, but the articles were not re-annotated by humans. Absolute scores on these labels therefore measure agreement with the categorization of publishers, especially for labels with unclear boundaries (\emph{society}, \emph{human interest}, \emph{lifestyle}).

MasakhaNEWS labels were assigned by human annotators under their own scheme and then mapped to IPTC. Only MN-DS is annotated directly with IPTC topics, and it is English only. The multilingual IPTC test set of Kuzman and Ljube\v{s}i\'c~\cite{kuzman2025llm} is not public, so it was not available for this study.

Label coverage per language is uneven, with a median of \val{ts.lang_labels.median} labels per language, and the African languages cover only four or five broad topics. Macro-F1 is computed over the labels present, so per-language scores are not directly comparable between languages with different label coverage. Publisher diversity is also limited for some languages: \val{ts.div.satu.langs} languages come from practically one publisher, namely BBC/VOA for the African languages, L3Cube for the Indian languages, and a single site for several CC-NEWS languages. Claims for these languages are limited to those publishers, and the site-cluster bootstrap widens their confidence intervals accordingly.

The teacher bounds the student from above, and its errors and biases are copied. The fidelity set measures similarity to the teacher, not correctness.

The hyperparameters of fastText and TF-IDF + SVM were set once without a search, whereas the StaticKD recipe was selected from \val{runs.total} runs on the development set. This unequal tuning budget may widen the accuracy gap to these two baselines. The classifier trainer of the Model2Vec library was not run as a baseline, and only its components were tested separately as ablations. MiniLM-L6 and DistilmBERT were trained and evaluated with at most 256 tokens and therefore see less than half of most test articles, so their accuracy may be underestimated.

The benchmark ran on one machine in one session, with \val{eff.teacher.n_latency} latency documents for the teacher and the transformers. Absolute numbers on a server will differ and the latency sample of the heavy models is small, but ratios between models measured under the same protocol remain informative.

StaticKD is much less accurate on African languages that are absent from both the distillation data and pre-training (Section~\ref{sec:discussion}). Used as a topic filter, it could misclassify or drop news in these languages more often. Before deployment, accuracy should be monitored per language, and low-confidence articles should be routed to the teacher through the cascade mode.

The study covers only top-level IPTC labels and single-label classification. Deeper levels and multi-topic articles remain untested.
